%%%%%%%%%%%%%%%%%%%%%%%%%%%%%%%%%%%%%%%%%%%%%%%%%%%%%%%%%%%%%%%%%%%%%%
% Lớp 11 | Quan hệ vuông góc trong không gian | Góc giữa đường thẳng và mặt phẳng. Góc nhị diện
% Dạng 1H8?6-4: Khoảng cách giữa điểm, đường, biết trước góc (𝑑, (𝑃))
% Số câu: 2  (sinh tự động bởi tach_cau.py; từ nay thầy sửa trực tiếp file này)
%%%%%%%%%%%%%%%%%%%%%%%%%%%%%%%%%%%%%%%%%%%%%%%%%%%%%%%%%%%%%%%%%%%%%%

%%%%%%%%%%%%%%%% TRẮC NGHIỆM 4 PHƯƠNG ÁN %%%%%%%%%%%%%%%%

% nguồn: Ok_TracNghiemLop11_Lan1
\begin{ex}%[1H8N6-4]%[ma:cc4bdaa0]
	Cho hình chóp tam giác đều $S.ABC$ có cạnh đáy bằng $2a$, góc giữa mặt bên và  mặt đáy bằng $60^\circ$. Độ dài cạnh bên của hình chóp bằng
	 \choice
		{\True $\dfrac{a\sqrt{21}}{3}$} 
		 {$\dfrac{a\sqrt{7}}{3}$}
		{$\dfrac{a\sqrt{17}}{3}$} 
		{$\dfrac{a\sqrt{15}}{3}$} 
	 \loigiai{
	\immini{Gọi $H$ là chân đường cao, suy ra $H$ là tâm tam  giác đều $ABC$.\\
		Gọi $M$ là trung điểm $BC$ suy ra góc giữa mặt bên  $(SBC)$ và $(ABC)$ là góc $\widehat{SMH}=60^\circ$.\\
		Ta  có $AH=\dfrac{2}{3}AM=\dfrac{2a\sqrt{3}}{3}; HM=\dfrac{1}{3}AM=\dfrac{a\sqrt{3}}{3}$.\\
		$SH=HM\cdot \tan  \widehat{SMH}=\dfrac{a\sqrt{3}}{3}\cdot \sqrt{3}=a$.\\
		Xét  $\triangle SHA: SA^2=SH^2+AH^2=\dfrac{4a^2}{3}+a^2=\dfrac{7a^2}{3}\Rightarrow SA=\dfrac{a\sqrt{21}}{3}$.}
		{\begin{tikzpicture}[>=stealth,line join=round,line cap=round,font=\footnotesize,scale=0.5]
		\clip (-1,-5)  rectangle (12,7);
		 \def\r{6.5};
		\def\g{120}; 
		\def\gC{30}; 
		\path 
			(0,0) coordinate  (A)
			(7,-4) coordinate  (B)
			(11,0) coordinate  (C)
			(6,6) coordinate  (S);
		\draw (S)--(A)--(B)--(C)--(S)--(B); 
		\coordinate (M) at  ($(B)!0.5!(C)$);
		\coordinate (H) at  ($(A)!0.666!(M)$);
		\draw[dashed]  (S)--(H)--(M)
			(C)--(A)--(H)--(B) 
			 (H)--(C)
			 (M)--(S);
		\fill[black] (A) circle (1pt) node[below left]{$A$} 
			(C) circle (1pt) node[below right]{$C$} 
			(B)  circle (1pt) node[below left]{$B$}
			 (S) circle (1pt) node[above right]{$S$}
			(H) circle  (1pt) node[below left]{$H$}
			(M) circle (1pt) node[below  right]{$M$};
		\draw  pic[draw,angle radius = 5pt] {right angle =S--H--B};
		\draw pic[draw,angle radius =  5pt] {right angle =S--H--C};
	\end{tikzpicture}
	 }}
\end{ex}


%%%%%%%%%%%%%%%% TRẢ LỜI NGẮN %%%%%%%%%%%%%%%%

% nguồn: Ok_TraLoiNgan_Lop11_Lan1
\begin{ex}%[1H8N6-4]%[ma:0fd4d60d]
	Cho hình chóp tam giác đều $S.ABC$ có cạnh đáy bằng $2\sqrt{21}$,  góc giữa mặt bên và mặt đáy bằng $60^\circ$. Tính độ dài cạnh bên của hình chóp.
	 \shortans{7}
	\loigiai{ 
	\immini{Đặt $a=\sqrt{21}$. Gọi $H$ là chân đường cao, suy ra $H$ là tâm  tam giác đều $ABC$.\\
		Gọi $M$ là trung điểm $BC$ suy ra góc giữa mặt bên $(SBC)$ và $(ABC)$ là  góc $\widehat{SMH}=60^\circ$.\\
		Ta có  $AH=\dfrac{2}{3}AM=\dfrac{2a\sqrt{3}}{3}; HM=\dfrac{1}{3}AM=\dfrac{a\sqrt{3}}{3}$.\\
		 $SH=HM\cdot \tan \widehat{SMH}=\dfrac{a\sqrt{3}}{3}\cdot \sqrt{3}=a$.\\
		Xét $\triangle SHA: SA^2=SH^2+AH^2=\dfrac{4a^2}{3}+a^2=\dfrac{7a^2}{3}\Rightarrow SA=\dfrac{a\sqrt{21}}{3}=7$.}{\begin{tikzpicture}[>=stealth,line join=round,line cap=round,font=\footnotesize,scale=0.5]
		\clip (-1,-5) rectangle  (12,8);
		\def\r{6.5}; 
		\def\g{120}; 
		 \def\gC{30};
		\path 
			 (0,0) coordinate (A)
			 (4,-4) coordinate (B)
			 (10,0) coordinate (C);
		\coordinate (M)  at ($(B)!0.5!(C)$);
		\coordinate (H) at ($(A)!0.666!(M)$); 
		 \coordinate (S') at ($(A)!(H)!(C)$);
		\coordinate (S) at  ($(H)!6!(S')$);
		 \draw (S)--(A)--(B)--(C)--(S)--(B)
			(M)--(S); 
		 \draw[dashed] (A)--(H)--(S)
			(A)--(C)--(H)--(M); 
		\fill[black]  (A) circle (1pt) node[above left]{$A$}
			(C) circle (1pt) node[below  right]{$C$}
			 (B) circle (1pt) node[below left]{$B$}
			(S) circle (1pt) node[above right]{$S$} 
			(M) circle (1pt) node[right]{$M$} 
			(H) circle (1pt)  node[below]{$H$};
		\draw pic[draw,angle radius  = 5pt] {right angle =S--H--A};
		\draw pic[draw,angle  radius = 5pt] {right angle =A--M--B};
	\end{tikzpicture}
	}} 
\end{ex}
